\section{ChronoLLM: domain-specific code assistance and the limits of simulation evidence}
\label{sec:exec24-chronollm}
A bank turning a collections question into executable analysis faces two
different tasks. It must express the intended decision, and it must implement
that decision in a changing software environment. Domain-specific code
assistance can reduce the second burden without solving the first. Wang and
coauthors study this distinction in an adjacent engineering setting:
ChronoLLM customizes language models to generate PyChrono scripts for
mechanical simulation. The inspected original is the 19 August 2025 arXiv
preprint, not the later journal version identified by the discovery index.
It supplies a useful design analogue for causal tools, rather than evidence
that its models identify bank interventions.\citep{exec24_chronollm_original}

\subsection{Version-bound knowledge and expert-curated training}
The source describes stale and mixed API knowledge, with examples of renamed
vector, gravity, collision and visualization interfaces. Its proposed training
pipeline combines continued pretraining on domain material with instruction
fine-tuning. The instruction collections include simulation generation,
natural-language-to-API mappings, API explanations and reasoning examples.
Expert contributors clean scripts and comments; model-generated descriptions
and transformations supplement that work. The paper also describes filtering,
privacy checks, deduplication and final quality review. These are reported
curation procedures, not evidence that we acquired or independently audited
the full training corpus.\citep{exec24_chronollm_original}

For the proposed causal library, this suggests binding every example to an
exact package version, method identifier and input/output contract. A
natural-language request for a default-risk comparison should not silently
select an obsolete method, change time zero, collapse competing events, or
replace a requested effect with a prediction. Retrieval and small curated
examples are plausible first prototypes; fine-tuning is a conditional
alternative when repeated version-specific failures justify its preparation
and maintenance cost. Porting the whole simulator or treating its generated
code as a causal engine would answer a different question.

The paper distinguishes in-context learning, which supplies examples without
changing model weights, from supervised fine-tuning and parameter-efficient
adaptation. Its broad architecture and model-selection discussion is not
adopted as a universal ranking. Our comparison would hold the selected
backend revision, prompt, package/API version, task family and allowed repair
budget fixed before estimating an improvement. An API spelling correction
and a correction to the scientific target must be separately observable.
\citep{exec24_chronollm_original}

\subsection{A score improvement depends on what the score measures}
The original distinguishes code similarity, execution tests and a separate
LLM judge. Its Table~7 reports similarity and judge scores under reference
code, documentation, or their combination. These are different measurements.
For the reported GPT-4o-mini profile, the combined reference/documentation
judge score rises from 41.80 to 68.03 after supervised fine-tuning, a difference
of 26.23 score units. The documentation-only score falls from 43.22 to 38.03,
a difference of $-5.19$. For the reported 8B profile, the combined judge score
is 39.06 with in-context learning and 38.78 with fine-tuning, a difference
of $-0.28$. Table~8 also contains lower judge scores for the 70B LoRA profile
than for its pretrained comparator. The exact subtractions reproduce the
displayed table entries; they are neither new experiments nor significance
tests.\citep{exec24_chronollm_original,exec24_chronollm_reference}

Thus the paper's broad improvement language needs metric-specific reading.
Similarity to reference code does not by itself establish functional or
physical correctness. A judge score is not a probability that the output
passes a causal identification contract. The text discusses pass-at-$k$ and
attributes judge-correlation findings to a companion study, but the inspected
Tables~6--8 do not supply a numerical pass-at-$k$ series. The conclusion's
mention of a Hacked SimBench baseline is not accompanied there by an
identified table row or reproducible comparison definition. These gaps call
for the companion original, exact task/split manifests, scorer revision and
run artifacts before borrowing a stronger empirical claim. They do not
justify silently inventing missing results.\citep{exec24_chronollm_original}

The two case studies show generated double-pendulum and vehicle scripts.
The authors acknowledge that their tasks may be in distribution relative to
training and describe the scripts as starting points for expert improvement.
The double-pendulum discussion distinguishes API errors from incorrect world
construction, axes and units. This is especially relevant to causal analysis:
a script can run while implementing the wrong population, assignment law or
intervention. The source reports example execution, but no acquired script,
simulator or model is executed in this proposal review.
\citep{exec24_chronollm_original}

\subsection{Executable worlds need not identify the intervention}
Consider an authored banking example with binary latent type $U$, equally
likely to be zero or one, and an observed action $A=U$. Let $Y$ denote
incremental account revenue in dollars. Two candidate structural worlds are
\begin{equation}
  \mathcal M_1:\ Y=10A+20U,
  \qquad
  \mathcal M_2:\ Y=30U.
  \label{eq:exec24-chronollm-worlds}
\end{equation}
Both generate the same observed pairs: $(A,Y)=(0,0)$ or $(1,30)$, each with
probability one half. The observed contrast is therefore USD~30 in either
world. Under an intervention that sets $A$ while preserving the distribution
of $U$, however, the average effect is USD~10 in $\mathcal M_1$ and zero in
$\mathcal M_2$. With a stipulated intervention cost of USD~5, the incremental
net value is USD~5 or USD~$-5$. The same observed data support opposite
decisions under these two structural explanations.
\citep{exec24_chronollm_reference}

Each world can be implemented with perfectly valid code and pass tests of its
own algebra. Execution cannot select which world describes the bank.
Conditional treatment overlap also fails here because action is fixed by
type. A useful agent should expose the unresolved assumption and request
appropriate additional information or a justified design. Public or generated
data remain valuable for testing whether the implementation preserves a
declared world and returns the required refusal; they do not turn that world
into an empirically verified bank mechanism. This provides a small exact
oracle for the proposed business-question-to-code workflow without paying
for an authentic bank experiment first.

\subsection{Full costs, repair accounting and conditional reuse}
The paper describes a small fine-tuning example taking roughly twenty minutes
and costing roughly USD~3 after the time-consuming data preparation was
complete. This is a reported historical example, not a current API quote,
full preparation cost, or measured bank productivity gain. Its model and task
profile must remain attached to the claim.\citep{exec24_chronollm_original}

For an independent stipulated accounting witness, combine an assumed USD~3
API charge with thirty minutes of expert work at USD~60 per hour. Total cost
is USD~33. Suppose eight requested tasks yield six compilable scripts, four
semantically valid simulations and two causally acceptable outputs. Cost is
USD~4.125 per request, USD~5.50 per compilable script and USD~16.50 per
accepted output. These invented inputs show why expert repair, refusal and
failure denominators must be retained; they are not a reconstruction of the
paper's experiment.\citep{exec24_chronollm_reference}

The proposed route is therefore reference use of its version-aware curation
and evaluation ideas, with an independently designed causal test suite.
The method owner defines intervention, units, clocks and admissible refusal;
the software owner pins package and interface behavior; the evaluation owner
holds out task families and versions, records all attempts and expert edits,
and compares retrieval, examples and fine-tuning under matched costs.
Unlearning stale APIs and richer tool or multi-agent interaction are future
directions in this preprint, not demonstrated universal remedies. Its
supplementary and declaration sections also use different availability
language: models/API access on request, a reported SimBench repository, and
a promised final code/data link. Exact artifacts, rights and reproducibility
therefore remain conditions for adoption.\citep{exec24_chronollm_original}

This adjacent study supports a concrete public-prototype design while leaving
scientific identification, native adoption and bank value unestablished.
